\usepackage{xcolor}
\usepackage{afterpage}
\usepackage{pifont,mdframed}
\usepackage[bottom]{footmisc}

\newcommand{\inputfile}{\texttt{stdin}}
\newcommand{\outputfile}{\texttt{stdout}}
\makeatletter
\renewcommand{\this@inputfilename}{\texttt{stdin}}
\renewcommand{\this@outputfilename}{\texttt{stdout}}
\makeatother

È l'ora di pranzo allo stage per le Olimpiadi di Informatica e tutti gli studenti sono affamati! Tommaso aveva promesso di preparare la pizza per tutti, ma come al solito è troppo lento. Samuele, stufo di aspettare, decide di andare lui stesso in pizzeria a prendere le pizze.

La città è composta da $N$ incroci, numerati da $1$ a $N$. Gli incroci sono collegati da $M$ strade bidirezionali. La $i$-esima strada collega gli incroci $a_i$ e $b_i$ e richiede $t_i$ minuti per essere percorsa. È sempre possibile raggiungere il Toppo partendo da un qualsiasi incrocio attraverso una sequenza di strade.

Ci sono $N$ partecipanti allo stage, uno per ogni incrocio da $1$ a $N$. Il Toppo (la sede dello stage dove tutti devono tornare) si trova all'incrocio $T$ ($1 \le T \le N$). Ogni partecipante deve tornare al Toppo partendo dal proprio incrocio.

Lungo il tragitto, in $K$ incroci ci sono dei chioschi che vendono tranci di pizza. Il chiosco nell'incrocio $i$ vende un trancio di pizza con un livello di soddisfazione $y_i$ (misurato in "quanto vale la pena fermarsi").

Ogni partecipante è disposto a fermarsi al massimo in un chiosco lungo il suo percorso verso il Toppo, ma solo se il tempo extra aggiunto al suo percorso è al massimo uguale al livello di soddisfazione della pizza che prenderebbe. In altre parole, se fermarsi a prendere una pizza gli fa perdere $t$ minuti in più rispetto al percorso più veloce, lo farà solo se il livello di soddisfazione è almeno $t$.

\begin{warning}
Un partecipante può passare attraverso incroci con chioschi senza fermarsi (semplicemente ignorando la pizza). Si ferma "ufficialmente" al massimo in un chiosco.
\end{warning}

\InputFile
La prima riga contiene quattro numeri interi $N$, $M$, $K$ e $T$ separati da spazi.

Ognuna delle successive $M$ righe contiene tre interi $a_i$, $b_i$ e $t_i$, che descrivono una strada tra gli incroci $a_i$ e $b_i$ che richiede $t_i$ minuti per essere percorsa ($a_i$ e $b_i$ sono diversi tra loro, e $t_i$ è un intero positivo al massimo $10^4$).

Le successive $K$ righe descrivono ciascuna un chiosco di pizza con due interi: l'indice del suo incrocio e il suo livello di soddisfazione (un intero positivo al massimo $10^9$). Più chioschi possono trovarsi nello stesso incrocio.

\OutputFile
L'output deve consistere di $N$ righe. La riga $i$ contiene il singolo intero $1$ se il partecipante nell'incrocio $i$ può fermarsi a mangiare una pizza lungo il percorso verso il Toppo, e $0$ altrimenti. (Anche il partecipante che si trova già al Toppo può decidere di muoversi per prendere una pizza e poi tornare, se il tempo extra lo permette).

\Constraints
\begin{itemize}
    \item $2 \le N \le 50\,000$.
    \item $1 \le M \le 100\,000$.
    \item $1 \le K \le N$.
    \item $1 \le T \le N$.
    \item $1 \le t_i \le 10^4$.
    \item $1 \le y_i \le 10^9$.
    \item Il grafo delle strade è connesso.
\end{itemize}

\Scoring
\begin{itemize}
    \item \subtask Casi d'esempio.
    \item \subtask $N \le 10$, $M \le 20$, $K \le 10$.
    \item \subtask $N \le 1\,000$, $M \le 2\,000$, $K \le 1\,000$.
    \item \subtask $K = 1$.
    \item \subtask $t_i = 1$ per ogni strada.
    \item \subtask Nessuna limitazione aggiuntiva.
\end{itemize}

\Examples
\begin{example}
    \exmpfile{input0.txt}{output0.txt}
\end{example}

\Explanations
Nel \textbf{caso d'esempio}:
\begin{itemize}
    \item Il partecipante nell'incrocio 3 dovrebbe fermarsi per una pizza: il suo percorso aumenterebbe solo di 6 minuti (da 2 a 8), e questo aumento è al massimo uguale al livello di soddisfazione 7 della pizza.
    \item Il partecipante nell'incrocio 2 dovrebbe ovviamente prendere la pizza nell'incrocio 2, dato che questo non causa alcun cambiamento nel suo percorso ottimale.
    \item Il partecipante nell'incrocio 1 è un caso interessante: potrebbe sembrare che il suo percorso ottimale (lunghezza 10) aumenterebbe troppo per giustificare una sosta per la pizza. Tuttavia, essendoci un chiosco all'incrocio 2 con una pizza molto soddisfacente, il partecipante decide di fare: $1 \to 4 \to 2$ (prende la pizza) $\to 4$ (Toppo). Questo percorso richiede un tempo extra di 6 minuti, che è accettabile poiché $6 \le 7$!
    \item Anche il partecipante nell'incrocio 4 (già al Toppo) deciderà di fare un salto all'incrocio 2: il suo percorso sarà $4 \to 2 \to 4$, con un tempo extra di 6 minuti. Poiché $6 \le 7$, deciderà di farlo.
\end{itemize}
